\cotefile{106}
\scanpage{2}
\margin{$\mathrm{Ho}\,C$\quad \unsure{Cat. d'obj. cofibrants}\quad ?}
Soit $(C,W_e,\mathrm{cof})$ une catégorie satisfaisant \unsure{C1 . C3|C1 : C3} de \uline{\unsure{Brown|Baues}}, (\unsure{pour mé. C0}).
\inserted{On suppose qu'il y a objet initial $\varnothing$}
(P. ex on pourrait avoir $\mathrm{cof}=\unsure{\mathrm{cof}_{W_e}}$). Alors il est connu que
\[
W_{\mathrm{Cof}}^{-1}\mathrm{Cof}\longrightarrow W^{-1}C
\]
est une équivalence de catégories.

Pour le voir, on va définir un foncteur en sens inverse (qui sera quasi-inverse)
\[
\Phi:W^{-1}C\longrightarrow W_{\mathrm{Cof}}^{-1}\mathrm{Cof},\quad\text{i.e.}\quad C\longrightarrow W_{\mathrm{Cof}}^{-1}\mathrm{Cof}
\]
qui rend les $u\in W$ invers.

Pour tt $x\in C$, utilisant C3 pour $\varnothing\longrightarrow x$, on trouve
\[
\tilde x\xrightarrow[\sim]{\alpha_x}x\qquad\text{avec}\quad \tilde x\in\mathrm{Cof},\ \alpha_x\in W.
\]
Soit alors $x\xrightarrow{f}y$ dans $C$,
\diagram{\begin{tikzcd}
\tilde x \arrow[d,"\alpha_x"',"\sim"] \arrow[dr,"f\alpha_x"] & \tilde y \arrow[d,"\alpha_y"]\\
x \arrow[r,"f"'] & y
\end{tikzcd}}
considérons
\diagram{\begin{tikzcd}
\unsure{\varnothing} \arrow[r] \arrow[d] & \tilde y \arrow[d,"\alpha_y","\sim"']\\
\tilde x \arrow[r,"f\alpha_x"'] & y
\end{tikzcd}}
appliquons le ``weak lifting lemma'', on trouve
\diagram{\begin{tikzcd}
& \tilde y \arrow[r,hook,"i","\sim"'] \arrow[d,"\unsure{\alpha_y}"',"\sim"] & \tilde y' \arrow[dl,"\unsure{\alpha'_y}","\sim"']\\
\tilde x \arrow[r,"f\alpha_x","g"'] \arrow[urr,bend left=12,"\tilde f"] & y &
\end{tikzcd}}
un relèvement $\tilde f$ \deleted{\unsure{comme un des}} \inserted{de $f\alpha_x$} \deleted{$f\alpha_x$} pas vers $\tilde y$, mais vers $\tilde y'$.

Considérons \unsure{lors|alors}
\[
[i]^{-1}\tilde f:\tilde x\longrightarrow\tilde y\qquad\text{dans }\mathrm{Hom}_{W_{\mathrm{Cof}}^{-1}}(\tilde x,\tilde y).
\]
Je dis qu'il ne dépend pas du choix de $(i,\alpha'_{\unsure{y}},\varphi)$. Pour deux tels choix, on
\diagram{\begin{tikzcd}
& \tilde y'_1\\
\tilde y \arrow[ur,hook,"i_1"] \arrow[dr,hook,"i_2"'] &\\
&\tilde y'_2
\end{tikzcd}}
\unsure{quitte à introduire} $\tilde y'_3=\tilde y'_1\vee_{\tilde y}\tilde y'_2$,
SPS que $\tilde y'_1=\tilde y'_2$, $i_1=i_2$, $\alpha'_1=\alpha'_2$, d'où \unsure{on a deux relèvements} $\varphi_1,\varphi_2$ de $f\alpha_x$ vers \unsure{$\underset{g=f\alpha_x}{y}=\tilde y'$}, \deleted{SPS $=$ \unsure{que $\tilde y$}} et on est ramené (écrivant $y'$ au lieu de $\tilde y'$, $x'$ au lieu de $\tilde x$, mais $x$ cof)
\diagram{\begin{tikzcd}
x' \arrow[r,shift left=.65ex,"\varphi"] \arrow[r,shift right=.65ex,"\psi"'] & y' \arrow[r,"\alpha","\sim"'] & y
\end{tikzcd}}
\diagram{\deleted{\begin{tikzcd}
& \unsure{y'}\ \mathrm{cof} \arrow[d,"\alpha","\sim"']\\
x'\ \mathrm{cof} \arrow[ur,shift left=.65ex,"\varphi"] \arrow[ur,shift right=.65ex,"\psi"'] \arrow[r,"g"'] & y
\end{tikzcd}}}
prouver que si $x',y'\in\mathrm{Cof}$, et $\alpha\varphi=\alpha\psi\deleted{\unsure{g}}$, alors
\[
[\varphi]=[\psi]\quad\text{dans } W_{\mathrm{Cof}}^{-1}\mathrm{Cof}.
\]
\margin{\diagram{\begin{tikzcd}
x' \arrow[r] & y\\
x'\amalg x' \arrow[u,"\nabla_{x'}"] \arrow[r,"(\varphi,\psi)"'] & y' \arrow[u,"\in W"']
\end{tikzcd}}
\diagram{\begin{tikzcd}
x'\ \mathrm{cof} \arrow[r,shift left=.65ex] \arrow[r,shift right=.65ex] & y'\ \mathrm{cof} \arrow[r,very thick] & y
\end{tikzcd}}}
\margin{\diagram{\begin{tikzcd}
\varnothing \arrow[r] & \unsure{x} \arrow[d]\\
& y
\end{tikzcd}}}
Si ce lemme est valable, on trouve un foncteur
\[
\deleted{\unsure{W}}C\longrightarrow W_{\mathrm{Cof}}^{-1}\mathrm{Cof},
\]
car la compatibilité de $u\mapsto\Phi(u)$ aux identités et aux compositions est immédiate.
\mspage{1}
